\begin{exo}{}
	On considère la fonction $f$ définie sur $\intfo{1}{+\infty}$ par :
	\[
		f(x) = \frac{12}{x}
	\]

	\begin{enumerate}
		\item Pourquoi ne pouvait-on pas définir la fonction $f$ sur $\intfo{0}{+\infty}$ ?
		\item Recopier et compléter, sans calculatrice, le tableau de valeurs suivant.
	\end{enumerate}

	\begin{center}
		\def\arraystretch{1.4}
		\begin{tabular}{|c|c|c|c|c|c|c|c|c|} \hline
			$x$    & $1$           & $2$           & $3$           & $4$           & $5$           & $6$           & $8$           & $10$          \\ \hline
			$f(x)$ & \phantom{xxx} & \phantom{xxx} & \phantom{xxx} & \phantom{xxx} & \phantom{xxx} & \phantom{xxx} & \phantom{xxx} & \phantom{xxx} \\ \hline
		\end{tabular}
	\end{center}

	\begin{enumerate}[resume]
		\item Tracer la courbe représentative de $f$ dans un repère orthonormé.
		\item Dresser son tableau de variations sur $\intff{1}{10}$ et préciser ses éventuels extremums.
	\end{enumerate}
\end{exo}
